\documentclass[10pt]{article}

\usepackage[utf8]{inputenc}

\usepackage[T1]{fontenc}

\usepackage{amsmath}

\usepackage{amsfonts}

\usepackage{amssymb}

\usepackage[version=4]{mhchem}

\usepackage{stmaryrd}



\begin{document}

формул сигнатуры $\langle\Omega, X\rangle$. Множество максимальных, совместных с $T(A, X)$, нножеств формул сигнатуры $\langle\Omega, X\rangle$, не содержащих свободных переменных, отличных от нек-рого фиксированного $v_{0}$, обозначается через $S(A, X)$. Теория $T$ наз. с т а б и л ь н о й в м о щ н о с т и $\lambda$, если для любой модели $A$ теории $T$ и любого $X \subseteq|A|$, мощность к-рого не превосходит $\lambda$, мощность $S(A, X)$ также не превосходит $\lambda$. Теория наз. с т а б и ли н о й, если она стабильна хотя бы в одной бесконечной мощности.



Пусть $|T|$ обозначает мощность множества формул сигнатуры $\Omega$. Если теория $T$ стабильна, то она стабильна во всех мощностях, удовлетворяющих равенству $\lambda=\lambda^{|T|}$. Если теория $T$ стабильна, то существуют модель $A$ теории $T$ и бесконечное множество $Y \subseteq|A|$ такие, что для любой формулы $\varphi\left(v_{1}, \ldots . v_{n}\right)$ сигнатуры $\Omega$ и для любых двух последовательностей $\left\langle a_{1}, \ldots, a_{n}\right\rangle$, $\left\langle b_{i}, \ldots, b_{n}\right\rangle$ различных элементов множества $Y$ истинность $\varphi\left(a_{1}, \ldots, a_{n}\right)$ в $A$ эквивалентна истинности $\varphi\left(b_{1}, \ldots, b_{n}\right)$ в $A$; при этом множество $Y$ наз. м н ож е с т в о м не р а з ли ч имы х в $T$ эл е мент о в. Оказывается, что характеристич. свойством н е с т а б и л ьны вание множества, имеющего в определенном смысле противоположные своӥства. А именно, нестабильность теории $T$ әквивалентна существованию формулы $\varphi\left(v_{1}\right.$, . ., $\left.v_{n} ; u_{1}, \ldots, u_{n}\right)$ сигнатуры $\Omega$, модели $A$ теории $T$ и последовательности $\left\langle a_{1}^{0}, \ldots, a_{n}^{0}\right\rangle, \ldots,\left\langle a_{1}^{k}, \ldots, a_{n}^{k}\right\rangle$, ... наборов элементов $A$ таких, что истинность $\varphi\left(a_{1}^{i}\right.$, $\left.\ldots, a_{n}^{i} ; a_{1}^{j}, \ldots, a_{n}^{j}\right)$ в $A$ равносильна неравенству $i<j$. Поэтому нестабильны полные расширения теории линейно упорядоченных множеств, имеющие бесконечные модели, а также теория любої бесконечной булевой алгебры. В частности, нестабильна теория натуральных чисел со сложением и теория поля действительных чисел. Если теория $T$ нестабильна, то число типов изоморфизма моделей $T$ в каждой несчетной мощности $\lambda>|T|$ равно $2^{\lambda}$. Поэтому теория $T$, категоричная в несчетной мощности $\lambda>|T|$, стабильна. Существуют, однако, стабильные теории, не категоричные ни в какой бесконечной мощности. Такова теория $T_{1}$, сигнатура к-рой состоит из одноместного предиката и счетного множества выделенных элементов. Аксиомы этой теории утверждают, что предикат истинен на выделенных элементах и делит каждую модель $T_{1}$ на два бесконечных множества, а также что выделенные элементы не равны между собой.



Теории конечной или счетной сигнатуры, стабильные в счетної мощности, наз. также т о т а л ь н о т р а н с ц е н д е н т ны м и. Всякая тотально трансцендентная теория стабильна во всех бесконечных мощностях. Всякая категоричная в несчетной мощности теория конечной или счетной сигнатуры является тотально трансцендентной. Упомянутая выше теория $T_{1}$ тотально трансцендентна. Тотально трансцендентные теории можно характеризовать и в других терминах. Пусть $T$ - полная теория конечной или счетной сигнатуры $\Omega, A$ - бесконечная модель теории $T$. Формуле $\varphi\left(v_{0}\right)$ сигнатуры $\langle\Omega,|A|\rangle$ припишем ранг -1 , если она ложна на всех элементах модели $\langle A,|A|\rangle$, и ранг $\alpha(\alpha-$ ординал), если ей не приписан никакой ранг, меньший $\alpha$, но для каждого әлементарного расширения $B$ системы $A$ и для каждой формулы $\psi\left(v_{0}\right)$ сигнатуры $\left\langle\Omega,|B|>\right.$ одной из формулі $\psi\left(v_{0}\right) \& \varphi\left(v_{0}\right)$ или $\rceil \psi\left(v_{0}\right) \& \varphi\left(v_{0}\right)$ приписан ранг, меньший $\alpha$. Теория $T$ тогда и только тогда является тотально трансцендентной, когда для каждой модели $A$ теории $T$ каждой формуле $\varphi$ сигнатуры $\langle\Omega,|A|\rangle$ приписан нек-рый ранг.



Jum.: [1] S h e la h S., "Ann. of math. logic", 1971, v. 3, № 3, p. 271-362; [2] e r o \% e, Classification theory and the number of non-isomorphic models, Amst.-[a. 0.], 1978 . $6 *$ E. А. Палютин, М. А. Тайчлин. СТАБИЛЬНЫЙ РАНГ к о л ь число $d$ такое, что любой $d$-порожденный проективный модуль над $R$ свободен. Кольцо $R$ здесь предполагается ассоциативно-коммутативным. Для некоммутативных колеп аналогичным образом определяемые левый и правый С. р. могут и не совпадать между собой.



СТАНДАРТИЗАЦИИ И УНИФИКАЦИИ МАТЕМА ТИЧЕСКИЕ ЗАДАЧИ - задачИ, в к-рых требуется определить оптимальные ряды изделий и их составных частей .



Оптимальшыӥ ряд изделий - это такоӥ набор различных типов изделий, взятых из исходного ряда, к-рый позволяет удовлетворить все заданные виды спроса в требуемом объеме с минимальными суммарными затратами на разработку, производство и әксплуатацию всех изделий. Оштимальный ряд изделий существует, поскольку с ростом числа типов изделий затраты на их разработку монотонно возрастают, а серийные и әксплуатационные затраты - убывают.



Терминологич. различие между задачами стандартизации и задачами униф্кикации является в определенноӥ степени условным. Существуют различные взгляды на вопрос о разграничении задач стандартизации и унификации. Напр., к задачам стандартизации относят задачи выбора оптимальных рядов относительно простых деталей и изделий, производимых массовыми сериями, тогда как к задачам униф্রикации относят задачи выбора оптимальных рядов сложных дорогостоящих изделий и их составных частей. Другой подход к разграничению задач стандартизации и унификации основан на учете степени детальности, с к-рой исследуется структура изделий, входящих в исходный ряд. Если изделия различного типа, входящие в исходныї ряд, полностью отличаотся друг от друга и не имеют одинаковых, т. е. унифицированных составных частей, то говорят об о дно у р о в н е в о й за д а ч е с т а нд а р т и з а ции, или просто о задаче с т ан д а рт и з а ц и и. Если изделия ряда рассматриваются с учетом их структуры и того, что изделия различного типа могут иметь унифицированные составные части, то говорят о двухуровневой задаче стандартизации. При дальнейшем рассмотрении структуры составных частей изделий можно получить $n$-уровневую задачу стандартизации. Задачи унификации - это $n$-уровневые задачи стандартизации с числом уровнеї $n>1$. Если предположить, что при определении оптимальных рядов сложных изделий, как правило, долюкны одновременно определяться оптимальные ряды их наиболее важных составных частей, то два описанных подхода к разграничению задач стандартизации и унификации совпадут.



Наиболее простым количественным методом, предназначенным для решения задач стандартизации при установлении рациональных параметров и размеров машин и оборудования, является использование системы предпочтительных чисел, основанной на применении геометрич. прогрессий. Установленные ряды предпочтительных чисел $R 5, \quad R 10, R 20, \quad R 40$ представляют собой ряды геометрич. прогрессий со знаменателями соответственно



\[

\begin{aligned}

\sqrt[5]{10} & \approx 1,6, \sqrt[10]{10} \approx 1,25, \\

\sqrt[20]{10} & \approx 1,12, \sqrt[40]{10} \approx 1,06 .

\end{aligned}

\]



Если для рассматриваемого класса изделий обоснована оптимальность одного из этих рядов и выбрано минимальное значение главного параметра $a_{0}$, то можно голучить значения главного параметра всех остальных изделиї ряда, округляя в случае необходимости величины $a_{0} q^{n}, n=1,2, \ldots$, где $q$ - знаменатель выбранного ряда.





\end{document}